% !TEX root = ../notes_missingsemester.tex
% Projektbericht Nachpruefung 2026: water rocket simulation tool.
%
% UNPUBLISHED. This file is deliberately NOT \input{} by
% notes_missingsemester.tex and its PDF is deliberately NOT written to
% chapter_pdfs/ (tools/update_readme_chapters.py generates the README table
% from that directory, and tools/build_chapter_pdfs.sh prunes anything there
% that the main file does not input). Build it standalone with:
%   tools/build_local_pdf.sh lecturenotes/notes_missingsemester 00_projektbericht_v2
% Output lands in lecturenotes/notes_missingsemester/local/.
%%%%%%%%%%%%%%%%%%%%%%%%%%%%%%%%%%%%%%%%%%%%%%%%%%%%%%%%%%%%%%%%%%%%%%

\index{Nachprüfung}\index{course work}
\chapter{Projektbericht Nachprüfung 2026}
\newthought{Real artists ship. - Steve Jobs}

\newthought{Build a water rocket simulation tool.}\index{water rocket}\index{simulation}\index{project report}
A user configures a rocket, bottle volume,
water fill level, launch pressure, nozzle diameter, fin and nosecone dimensions and position, masses, and then flies it in
simulation. The primary quantity is the apogee. Make sure to survey what already exists in the state of the art and build from there.

\marginnote{\newthought{How to write the report.} A practical guide for writing reports and theses lives \href{https://pages.schutera.com/}{here}.

\medskip
\null\hfill\qrcode[height=1.0cm]{https://pages.schutera.com/}
\medskip

For deeper druidic knowledge on writing a thesis, consult \href{https://pages.schutera.com/HowToThesis/notes.pdf}{here}.

\medskip
\null\hfill\qrcode[height=1.0cm]{https://pages.schutera.com/HowToThesis/notes.pdf}
\medskip}

\newthought{The tool is a full stack application}\index{full stack}: database, backend and
frontend as separate services, wired along the practices of the lecture blocks, and deployable
on any machine in one step with what you learned in the lecture, a docker compose or an equivalent.
The objective is a properly designed simulation tool, not the optimal rocket configuration, so results reports the evidence that yours behaves as
specified: A validation against a reference tool or an analytical estimate, or your own measurements
against predefined dimensions and KPIs.

\newthought{At least one algorithm from Fortgeschrittene Programmierung}\index{algorithm}
must be implemented in the tool and carry a functional part of it. Reference the script of the lecture block and the original paper, and explain how it is used in your tool. 

\newthought{Submit a written report as PDF}\index{assessment} and a link to your GitHub
repository. Seven pages in IMRD structure\index{IMRD}: five on the tool and the full stack,
two on the algorithms taken from the other modules.

\marginnote{\newthought{Seven pages, and how they split.}
Five pages on the tool and the full stack: problem, state of the art, design, architecture, data model,
deployment, results, and discussion including limitations and next steps.
Two pages on the algorithms from Fortgeschrittene Programmierung, the algorithm may run across Introduction, Methods, Results and Discussion.}


\newthought{Evaluation} weighs the scope of the State of the Art Research, accuracy of the simulated apogee, the realism and range of
the rocket configuration, the quality of the data model, the ease of large scale experimentation,
the depth of the algorithmic work, and the quality of the writing and structure.
